Our primary empirical evaluation tests whether the proposed spectrum-aware
parameterization remains effective on standard natural language understanding
tasks under compact trainable parameter budgets. We evaluate UILinLoRA and
UIOrthoLoRA on six GLUE validation tasks using RoBERTa-base and RoBERTa-large. Full training
details, adapter placement, optimization settings, and evaluation metrics are
provided in Appendix~\ref{app:training_details}.

We compare against representative PEFT baseline results reported by
\citet{albert2025randlora}, including VeRA, LoRA, and RandLoRA, and report
validation performance together with the corresponding trainable parameter
budgets. As shown in 
Table~\ref{tab:glue_results_base_large}, UIOrthoLoRA achieves the strongest
RoBERTa-base average among the compared methods, reaching 85.5 with 0.4M
trainable parameters. UILinLoRA, which removes the orthogonal rotations and
keeps only the magnitude-based spectral adaptation with diagonal scalers,
remains highly competitive, reaching 84.68 with only 0.08M trainable parameters
on RoBERTa-base and achieving the strongest RoBERTa-large average among the
compared proposed variants. These results suggest that pretrained singular
coordinates provide an effective basis for parameter-efficient adaptation, while
compact orthogonal rotations can provide additional expressiveness in some
settings without substantially increasing the trainable footprint.

Beyond the main GLUE results, we include three additional empirical studies in
the appendix. Appendix~\ref{app:soft_mixing_ablation} reports a diagnostic
soft-mixing regularization ablation that directly tests the leading-tail
mixing quantities from our analysis. Appendix~\ref{app:gpt2_generation_results}
reports GPT-2 Medium generation results, and
Appendices~\ref{app:retention_adaptation_protocol}
through~\ref{app:retention_adaptation_results} report a retention-oriented
evaluation protocol and results, which further motivate the
preservation-adaptation perspective underlying UIOrthoLoRA.

\subsection{Soft-Mixing Regularization Ablation}
\label{app:soft_mixing_ablation}

This section evaluates whether the scaler-induced mixing quantities derived in
Section~\ref{sec:uiortholora} can be used as practical regularizers. The goal is
not to introduce a new primary training recipe, but to provide a diagnostic test
of the preservation mechanism suggested by the theory.

For the UI update
\[
\Delta W^{ui}
=
E \bar{U}_r \Sigma'_r \bar{V}_r^\top D,
\]
our analysis identifies the left and right leading-tail mixing terms
\[
M_E = U_R^\top E \bar{U}_r,
\qquad
M_D = V_R^\top D \bar{V}_r.
\]
We add the soft-mixing penalty
\[
\mathcal{R}_{\mathrm{mix}}
=
\lambda_E^{\mathrm{mix}}\|M_E\|_F^2
+
\lambda_D^{\mathrm{mix}}\|M_D\|_F^2,
\]
and train with
\[
\mathcal{L}_{\mathrm{total}}
=
\mathcal{L}_{\mathrm{task}}
+
\mathcal{R}_{\mathrm{mix}}.
\]

\paragraph{Setup.}
We run this ablation on five GLUE tasks, RTE, MRPC, CoLA, STS-B, and
SST-2, using RoBERTa-base. The training recipe, adapter placement,
optimizer, scheduler, and task hyperparameters are kept identical to the
corresponding GLUE runs. Adapters are inserted into the query, key, value, and
attention-output projections in all 12 transformer layers. For each adapted
matrix, we adapt 256 singular values and set \(k_{vec}=0\), so no orthogonal
tail rotations are used in this diagnostic experiment. All reported values are averaged over two seeds. Spectral metrics are averaged over the 48 adapted attention
modules for each seed and then averaged across seeds.

We compare three variants:
\[
\text{A: } \lambda_E^{\mathrm{mix}}=0,\quad \lambda_D^{\mathrm{mix}}=0,
\]
\[
\text{B: } \lambda_E^{\mathrm{mix}}=10^{-3},\quad
\lambda_D^{\mathrm{mix}}=0,
\]
\[
\text{C: } \lambda_E^{\mathrm{mix}}=10^{-3},\quad
\lambda_D^{\mathrm{mix}}=10^{-3}.
\]
Variant A is the unregularized adapter, variant B penalizes only left-side
mixing, and variant C penalizes both left- and right-side mixing.

\paragraph{Metrics.}
For each adapted layer, we compute the mixing magnitudes
\[
\mu_E = \|M_E\|_2,
\qquad
\nu_D = \|M_D\|_2,
\]
the relative update magnitude
\[
\mathrm{RelPert}_F =
\frac{\|\Delta W^{ui}\|_F}{\|W^{pre}\|_F},
\]
and the update in pretrained singular coordinates
\[
C = U^\top \Delta W^{ui} V
=
\begin{bmatrix}
C_{11} & C_{12}\\
C_{21} & C_{22}
\end{bmatrix}.
\]
We report the leading-block leakage
\[
\mathrm{Leak}_{11}
=
\frac{\|C_{11}\|_2}{\|\Delta W^{ui}\|_2},
\]
and the off-tail ratio
\[
\mathrm{OffTailRatio}_F
=
\frac{\|C_{\mathrm{off}}\|_F}{\|C\|_F},
\qquad
C_{\mathrm{off}}
=
\begin{bmatrix}
C_{11} & C_{12}\\
C_{21} & 0
\end{bmatrix}.
\]
We also report leading left and right subspace drift using the operator-norm
\(\sin\Theta\) distance between the leading singular subspaces of
\(W^{pre}\) and \(W^{pre}+\Delta W\).

\begin{table}[h]
\centering
\caption{Soft-mixing regularization ablation on five GLUE tasks. Values are
averaged over two seeds. Spectral metrics are averaged across the 48 adapted
attention modules for each seed and then averaged across seeds. Variants: A is
no regularization, B is the \(\mu_E\)-only penalty, and C is the
\(\mu_E+\nu_D\) penalty. RTE and SST-2 report accuracy, MRPC reports F1, CoLA
reports Matthews correlation, and STS-B reports Pearson correlation.}
\label{tab:soft_mixing_ablation}
\begin{adjustbox}{max width=\linewidth}
\begin{tabular}{llccccccccc}
\toprule
\textbf{Task} &
\textbf{Variant} &
\(\lambda_E^{\mathrm{mix}}\) &
\(\lambda_D^{\mathrm{mix}}\) &
\textbf{Score} &
\(\mu_E\) &
\(\nu_D\) &
\(\mathrm{RelPert}_F\) &
\(\mathrm{Leak}_{11}\) &
\(\mathrm{OffTailRatio}_F\) &
\(\mathrm{Drift}_U / \mathrm{Drift}_V\) \\
\midrule
RTE   & A & 0 & 0 & 0.7708 & 1.384 & 1.442 & 0.379 & 0.722 & 0.948 & 1.000 / 1.000 \\
RTE   & B & \(10^{-3}\) & 0 & 0.7906 & 0.278 & 2.112 & 0.285 & 0.633 & 0.934 & 0.999 / 0.999 \\
RTE   & C & \(10^{-3}\) & \(10^{-3}\) & 0.7672 & 0.221 & 0.204 & 0.063 & 0.590 & 0.838 & 0.523 / 0.512 \\
\midrule
MRPC  & A & 0 & 0 & 0.9215 & 1.282 & 1.345 & 0.318 & 0.722 & 0.948 & 0.999 / 1.000 \\
MRPC  & B & \(10^{-3}\) & 0 & 0.9240 & 0.259 & 1.538 & 0.196 & 0.697 & 0.948 & 0.983 / 0.983 \\
MRPC  & C & \(10^{-3}\) & \(10^{-3}\) & 0.9132 & 0.230 & 0.188 & 0.043 & 0.626 & 0.857 & 0.667 / 0.662 \\
\midrule
CoLA  & A & 0 & 0 & 0.6409 & 1.249 & 1.336 & 0.308 & 0.721 & 0.946 & 0.999 / 0.999 \\
CoLA  & B & \(10^{-3}\) & 0 & 0.6429 & 0.287 & 1.618 & 0.206 & 0.637 & 0.938 & 0.989 / 0.989 \\
CoLA  & C & \(10^{-3}\) & \(10^{-3}\) & 0.6315 & 0.160 & 0.144 & 0.034 & 0.580 & 0.880 & 0.559 / 0.552 \\
\midrule
STS-B & A & 0 & 0 & 0.9080 & 1.375 & 1.447 & 0.382 & 0.724 & 0.947 & 1.000 / 1.000 \\
STS-B & B & \(10^{-3}\) & 0 & 0.9048 & 0.275 & 1.722 & 0.251 & 0.667 & 0.940 & 1.000 / 1.000 \\
STS-B & C & \(10^{-3}\) & \(10^{-3}\) & 0.9058 & 0.202 & 0.248 & 0.068 & 0.508 & 0.806 & 0.749 / 0.746 \\
\midrule
SST-2 & A & 0 & 0 & 0.9472 & 1.219 & 1.255 & 0.282 & 0.731 & 0.948 & 1.000 / 1.000 \\
SST-2 & B & \(10^{-3}\) & 0 & 0.9438 & 0.185 & 1.604 & 0.268 & 0.707 & 0.943 & 1.000 / 1.000 \\
SST-2 & C & \(10^{-3}\) & \(10^{-3}\) & 0.9369 & 0.134 & 0.170 & 0.056 & 0.565 & 0.850 & 0.685 / 0.672 \\
\bottomrule
\end{tabular}
\end{adjustbox}
\end{table}

\paragraph{Results.}
Across all five tasks, penalizing both mixing terms reduces the measured
interaction between the adapted tail components and the preserved leading
singular structure. Variant C substantially reduces both \(\mu_E\) and
\(\nu_D\), while variant B reduces \(\mu_E\) but leaves \(\nu_D\) large. This
one-sided behavior is consistent across the seed-averaged results: when only
\(\lambda_E^{\mathrm{mix}}\) is active, \(\nu_D\) increases from variant A to
variant B on every task. Thus, penalizing only one side allows the model to
route mixing through the unpenalized right-side scaler. The full regularizer
avoids this compensation, reducing \(\nu_D\) to values between \(0.144\) and
\(0.248\).

The leakage metrics show the same trend. Moving from the unregularized adapter
to the full regularizer reduces the mean off-tail ratio on every task. The
decrease is \(0.948 \rightarrow 0.838\) on RTE,
\(0.948 \rightarrow 0.857\) on MRPC,
\(0.946 \rightarrow 0.880\) on CoLA,
\(0.947 \rightarrow 0.806\) on STS-B, and
\(0.948 \rightarrow 0.850\) on SST-2. The leading-block leakage
\(\mathrm{Leak}_{11}\) also decreases under the full regularizer on every task.

The leading-subspace drift metrics provide the clearest separation between the
one-sided and two-sided penalties. Variants A and B remain close to maximal
mean drift on most tasks, while variant C consistently reduces the drift of the
leading singular subspaces. The mean left/right drift decreases from
approximately \(1.0/1.0\) under the unregularized adapter to \(0.523/0.512\)
on RTE, \(0.667/0.662\) on MRPC, \(0.559/0.552\) on CoLA, \(0.749/0.746\) on
STS-B, and \(0.685/0.672\) on SST-2. Since this drift is measured with an
operator-norm \(\sin\Theta\) distance, values near one indicate that at least
one leading direction undergoes a large angular change, rather than implying
that the entire leading subspace is orthogonal. Nevertheless, the consistent
reduction under variant C indicates improved preservation of the leading
spectral geometry.

\paragraph{Task performance and update magnitude.}
The full regularizer keeps downstream performance close to the unregularized
adapter while substantially reducing spectral leakage. The task-level score
decreases from variant A to variant C are \(0.36\) points on RTE, \(0.84\)
points on MRPC, \(0.95\) points on CoLA, \(0.22\) points on STS-B, and
\(1.03\) points on SST-2. Thus, the full regularizer produces a clear
structural preservation effect while keeping task performance close to the
unregularized adapter, with deviations of roughly one point across the
evaluated settings.

The relative update magnitude decreases under the full regularizer, but it
does not vanish. Mean \(\mathrm{RelPert}_F\) under variant C remains nonzero
across all tasks, with values in the approximate range of \(0.03\) to \(0.07\).
At the seed-resolved level, the reduction from variant A to variant C ranges
from roughly \(5\times\) to \(17\times\), while the off-tail and drift metrics
continue to improve. This indicates that the regularizer is not merely
shrinking the update norm; it also changes where the update lies in the
pretrained singular coordinate system.

Variant B obtains the strongest downstream score on RTE, MRPC, and CoLA, but
it does so while leaving substantial right-side mixing. This supports the
interpretation that partial mixing control can act as an adaptation bias, but
not as a sufficient preservation mechanism. Variant C is therefore the more
diagnostic setting for testing the preservation claim: it is the only variant
that consistently suppresses both sides of scaler-induced mixing and produces
the strongest reduction in off-tail leakage and leading subspace drift.

\paragraph{Takeaway.}
This diagnostic ablation supports the interpretation of \(\mu_E\) and
\(\nu_D\) as meaningful quantities for controlling scaler-induced leakage.
Penalizing only \(U_R^\top E\bar{U}_r\) reduces left-side mixing but leaves a
consistent escape route through \(V_R^\top D\bar{V}_r\). Penalizing both
\(U_R^\top E\bar{U}_r\) and \(V_R^\top D\bar{V}_r\) gives the cleanest
preservation behavior: it reduces both mixing magnitudes, lowers off-tail
leakage, reduces leading-block contamination, and substantially decreases
leading-subspace drift across tasks and seeds, while keeping downstream
performance close to the unregularized adapter.

\paragraph{Future directions.}
The present ablation suggests that the mixing magnitudes \(\mu_E\) and
\(\nu_D\) are not only useful analytical quantities, but may also serve as
practical control signals during training. Since reducing these quantities also
reduces off-tail leakage and leading-subspace drift, an interesting direction
is to study whether such spectral control can improve knowledge preservation
under broader retention-oriented adaptation settings.

In particular, future work may investigate whether controlling scaler-induced
mixing helps limit functional forgetting when adapters are trained to acquire
new task-specific knowledge. This would connect the spectral preservation
metrics studied here with the retention-adaptation perspective described in
Appendix~\ref{app:retention_adaptation_protocol}. Overall, this suggests a
possible path toward PEFT methods that regulate not only how much they adapt,
but also how much pretrained structure is allowed to drift during fine-tuning.
